% 取消下一行注释可生成 handout（无逐步显示）；保持注释则生成正常 slide。
% \PassOptionsToClass{handout}{beamer}
\documentclass[Prob-All.tex]{subfiles}
\begin{document}
\title[概率论]{第十四讲：条件分布与随机变量的独立性}
%\author[张鑫 {\rm Email: x.zhang.seu@foxmail.com} ]{\large 张 鑫}
\institute{\large \textrm{Email: xzhangseu@seu.edu.cn} \\ \quad  \\
	\vspace{0.3cm}
%	\trc{公共邮箱: \textrm{zy.prob@qq.com}\\
	%	\hspace{-1.7cm}  密 \qquad 码: \textrm{seu!prob}}
}
\date{}





{ \setbeamertemplate{footline}{}
	\begin{frame}
		\titlepage
	\end{frame}
}
%\section{随机变量}
\subsection{条件分布}
\begin{frame}
	\frametitle{条件分布}
	\begin{itemize}[<+-|alert@+>]
		\item $(\Omega,\mathcal{F},P)$ 为一概率空间，任意固定 $B\in\mathcal{F}$ 且 $P (B)>0$, $P (\cdot|B)$ 仍为 $(\Omega,\mathcal{F})$ 上的概率；
		\item 对于 $(\Omega,\mathcal{F})$ 上的随机变量 $X$, 考虑其关于条件概率 $P (\cdot|B)$ 的分布即在概率空间 $(\Omega,\mathcal{F},P (\cdot|B))$ 上考虑 $X$ 的分布:
		\begin{eqnarray*}
			F(x|B):=P(X\le x|B)=\dfrac{P(X\le x, B)}{P(B)}, x\in R
		\end{eqnarray*}
		\item 现假定有另一随机变量 $Y$, 则当 $P (Y=y)>0$ 时，$F (x|Y=y)$ 则为已知 $Y=y$ 时 $X$ 的条件分布函数.
	\end{itemize}

\end{frame}

\begin{frame}
	\frametitle{离散型条件分布}
	\begin{defi}
		若离散型随机向量 $(X,Y)$ 的联合分布为 $\{p_{ij}\}$, 则当 $$P (Y=y_j)=p_{\cdot j}>0$$ 时，称
		\begin{eqnarray*}
			p_{i|j}:=P(X=x_i|Y=y_j)=\dfrac{p_{ij}}{p_{\cdot j}}, i=1,2,\cdots,
		\end{eqnarray*}
		为已知 $Y=y_j$ 时 $X$ 的条件分布. \pause 类似的，当 $P (X=x_i)=p_{i\cdot}>0$ 时，称
		\begin{eqnarray*}
			p_{j|i}:=P(Y=y_j|X=x_i)=\dfrac{p_{ij}}{p_{i\cdot}}
		\end{eqnarray*}
		为已知 $X=x_i$ 时 $Y$ 的条件分布.
	\end{defi}

\end{frame}
\begin{frame}
	\frametitle{离散型条件分布函数}
	根据上面的定义，给定 $Y=y_j$ 条件下 $X$ 的条件分布函数为
	\begin{eqnarray*}
		F(x|y_j)=\sum_{i:x_i\le x}P(X=x_i|Y=y_j)=\sum_{i:x_i\le x}p_{i|j}
	\end{eqnarray*}
	\pause 类似的给定 $X=x_i$ 条件下 $Y$ 的条件分布函数为
	\begin{eqnarray*}
		F(y|x_i)=\sum_{j:y_j\le y}P(Y=y_j|X=x_i)=\sum_{j:y_j\le y}p_{j|i}
	\end{eqnarray*}

\end{frame}
\begin{frame}
	\frametitle{连续型条件分布}
	\begin{itemize}[<+-|alert@+>]
		\item 设二维连续型随机变量 $(X,Y)$ 的联合密度函数为 $p (x,y)$, 边际密度函数分别为 $p_X (x), p_Y (y)$;
		\item 条件分布函数 $P (X\le x|Y=y)$: 因为 $P (Y=y)=0$, 故无法通过条件概率直接计算；
		\item $P(X\le x|Y=y):=\lim_{h\rightarrow 0}P(X\le x|Y\in[y,y+h])$;
		\item 根据如上定义:\vspace{-0.3cm}
		{\small \begin{eqnarray*}
				\hspace{-0.5cm} F(x|y)&:=&P(X\le x|Y=y)=\pause \lim_{h\rightarrow 0}\dfrac{P(X\le x, Y\in[y,y+h])}{P(Y\in[y,y+h])}\\
				&=&\pause \lim_{h\rightarrow 0}\dfrac{\int_{-\infty}^x\int_y^{y+h}p(u,v)dvdu}{\int_y^{y+h}p_Y(v)dv}
				=\pause \lim_{h\rightarrow 0}\dfrac{\int_{-\infty}^x\big\{\dfrac{1}{h}\int_y^{y+h}p(u,v)dv\big\}du}{\dfrac{1}{h}\int_y^{y+h}p_Y(v)dv}\\
				&&\hspace{-0.7cm}\pause \xlongequal[\mbox{积分中值定理}]{p (x,y),p_Y (y)\mbox{在} y\mbox{连续}}\pause \int_{-\infty}^x\dfrac{p (u,y)}{p_Y (y)} du\pause
		\end{eqnarray*}}%
		\item 这表明给定 $Y=y$ 时 $X$ 的条件分布仍为连续型，其密度函数为 $$p (x|y)=p (x,y)/p_Y (y).$$
	\end{itemize}
\end{frame}
\begin{frame}
	\frametitle{连续随机变量的条件分布函数及条件密度函数}
	\begin{defi}
		对一切使 $p_Y (y)>0$ 的 $y$, 给定 $Y=y$ 条件下 $X$ 的条件分布函数和条件密度函数分别为
		\begin{eqnarray*}
			F(x|y)&=&\int_{-\infty}^x\dfrac{p(u,y)}{p_Y(y)}du,\\
			p(x|y)&=&\dfrac{p(x,y)}{p_Y(y)}.
		\end{eqnarray*}
		同理，对一切使 $p_X (x)>0$ 的 $x$, 给定 $X=x$ 条件下 $X$ 的条件分布函数和条件密度函数分别为
		\begin{eqnarray*}
			F(y|x)&=&\int_{-\infty}^y\dfrac{p(x,v)}{p_X(x)}dv,\\
			p(y|x)&=&\dfrac{p(x,y)}{p_X(x)}.
		\end{eqnarray*}
	\end{defi}
\end{frame}
\begin{frame}
	\frametitle{连续场合的全概率公式和贝叶斯公式}
	\begin{itemize}[<+-|alert@+>]
		\item 由条件密度函数的定义可得
		\begin{eqnarray*}
			p(x,y)=p_X(x)p(y|x), \quad p(x,y)=p_Y(y)p(x|y)
		\end{eqnarray*}
		\item 对 $p (x,y)$ 求边际密度函数可得
		\begin{eqnarray*}
			p_X(x)&=&\pause \int_{-\infty}^{+\infty}p_Y(y)p(x|y)dy\pause \\
			\pause P(X\le x)&=&\pause \int_{-\infty}^x\bigg[\int_{-\infty}^{+\infty}p_Y(y)p(u|y)dy\bigg]du\\
			&=&\pause \int_{-\infty}^{+\infty}\bigg[\int_{-\infty}^xp(u|y)du\bigg]p_Y(y)dy\\
			&=&\pause \int_{-\infty}^{+\infty}P(X\le x|Y=y)dP(Y\le y)
			% p_Y(y)=\int_{-\infty}^{+\infty}p_X(x)p(y|x)dx, F(y)=\int_{-\infty}^{+\infty}\bigg[\int_{-\infty}^yp(v|x)dv\bigg]p_X(x)dx\\
		\end{eqnarray*}

	\end{itemize}
\end{frame}
\begin{frame}
	\frametitle{连续场合的全概率公式和贝叶斯公式}
	\begin{itemize}[<+-|alert@+>]
		\item 类似推导可得
		\begin{eqnarray*}
			p_Y(y)&=&\pause \int_{-\infty}^{+\infty}p_X(x)p(y|x)dx \\
			\pause P(Y\le y)&=& \pause \int_{-\infty}^{+\infty}P(Y\le y|X=x)dP(X\le x)
			% p_Y(y)=\int_{-\infty}^{+\infty}p_X(x)p(y|x)dx, F(y)=\int_{-\infty}^{+\infty}\bigg[\int_{-\infty}^yp(v|x)dv\bigg]p_X(x)dx\\
		\end{eqnarray*}
		\item 贝叶斯公式的密度函数形式:
		\begin{eqnarray*}
			p(x|y)=\dfrac{p(x,y)}{p_Y(y)}=\dfrac{p_X(x)p(y|x)}{\int_{-\infty}^{+\infty}p_X(x)p(y|x)dx}\\
			p(y|x)=\dfrac{p(x,y)}{p_X(x)}=\dfrac{p_Y(y)p(x|y)}{\int_{-\infty}^{+\infty}p_Y(y)p(x|y)dy}
		\end{eqnarray*}

	\end{itemize}
\end{frame}


\begin{frame}
	\begin{exam}
		若 $(X,Y)$ 为二维正态分布 $N (\mu_1,\mu_2,\sigma_1^2,\sigma_2^2,\rho)$, 试求给定 $Y=y$ 时 $X$ 的条件分布密度.\\
	\end{exam}


	\pause\jieda $X$ 的条件分布密度为
	\begin{eqnarray*}
		&&\hspace{-0.8cm}p(x|y)=\pause \dfrac{p(x,y)}{p_Y(y)}\\
		&=&\pause  \dfrac{\frac{1}{2\pi \sigma_1\sigma_2\sqrt{1-\rho^2}}\exp\big\{-\frac{1}{2(1-\rho^2)}\big[\frac{(x-\mu_1)^2}{\sigma_1^2}-\frac{2\rho(x-\mu_1)(y-\mu_2)}{\sigma_1\sigma_2}+\frac{(y-\mu_2)^2}{\sigma_2^2}\big]\big\}}{\frac{1}{\sqrt{2\pi}\sigma_2}\exp\big\{-\frac{(y-\mu_2)^2}{2\sigma_2^2}\big\}}\\
		&=&\pause \frac{\exp\big\{\frac{-1}{2\sigma_1^2(1-\rho^2)}\big[(x-\mu_1)^2-2\rho\frac{\sigma_1}{\sigma_2}(x-\mu_1)(y-\mu_2)+\rho^2\frac{\sigma_1^2}{\sigma_2^2}(y-\mu_2)^2\big]\big\}}{\sqrt{2\pi} \sigma_1\sqrt{1-\rho^2}}\\
		&=&\pause \frac{1}{\sqrt{2\pi} \sigma_1\sqrt{1-\rho^2}}\exp\big\{-\frac{1}{2\sigma_1^2(1-\rho^2)}\big[x-(\mu_1+\rho\frac{\sigma_1}{\sigma_2}(y-\mu_2)\big]^2\big\}\\
		&\sim&\pause  N(\mu_1+\rho\frac{\sigma_1}{\sigma_2}(y-\mu_2),\sigma_1^2(1-\rho^2)) \xlongequal{\rho=0} N(\mu_1,\sigma_1^2)
	\end{eqnarray*}
	\pause 类似可求得 $Y$ 关于 $X=x$ 的条件分布密度.
\end{frame}

\begin{frame}
	\begin{exam}
		设 $(X,Y)$ 服从单位圆 $D=\{(x,y):x^2+y^2<1\}$ 上的均匀分布，试求条件分布密度函数 $p (y|x)$.
	\end{exam}

	\jieda
	$(X,Y)$ 的联合密度及 $X$ 的边缘密度分别为
	\begin{eqnarray*}
		p(x,y)=\left\{
		\begin{array}{ll}
			\dfrac{1}{\pi}, & x^2+y^2<1,\\
			0, & \mbox{其他}
		\end{array}\right. ,\quad
		p_X(x)=\left\{\begin{array}{ll}
			\dfrac{2}{\pi}\sqrt{1-x^2}, & |x|<1,\\
			0, & \mbox{其他}.
		\end{array}\right.
	\end{eqnarray*}
	\pause 故
	\begin{eqnarray*}
		p(y|x)=\dfrac{p(x,y)}{p_X(x)}=\left\{
		\begin{array}{ll}
			\dfrac{1}{2\sqrt{1-x^2}}, & |y|\le \sqrt{1-x^2}\\
			0, &\mbox{其他}
		\end{array}
		\right.
	\end{eqnarray*}

\end{frame}


\subsection{随机变量的独立性}

\begin{frame}
	\frametitle{随机变量的独立性}
	\begin{itemize}[<+-|alert@+>]
		\item 首先我们回顾一下事件的独立性: $$P (AB)=P (A) P (B)\Leftrightarrow P (A|B)=P (A), \mbox{当} P (B)>0;$$
		\item 二维离散型随机变量
		\begin{eqnarray*}
			p_{i|j}:=\textcolor{red}{\dfrac{p_{ij}}{p_{\cdot j}}}=P (X=x_i|Y=y_j)\pause\xlongequal[\mbox{不影响概率}]{\mbox{若条件事件}} P (X=x_i)\textcolor{red}{=p_{i\cdot}};
		\end{eqnarray*}
		\item 故如果对任意的事件 $\{X=x_i\}, \{Y=y_j\}, i,j=1,2,\cdots,$ 均独立，我们则可称随机变量 $X,Y$ 独立，即 $X,Y$ 独立等价于
		\begin{eqnarray*}
			p_{ij}=p_{i\cdot}p_{\cdot j}, \quad i,j=1,2,\cdots,
		\end{eqnarray*}
		\item 此时的联合分布函数 $F (x,y)$ 为
		\begin{eqnarray*}
			\textcolor{red}{F(x,y)}&=&\pause P(X\le x,Y\le y)=\pause \sum_{x_i\le x}\sum_{y_j\le y}p_{ij}=\pause \sum_{x_i\le x}p_{i\cdot}\sum_{y_j\le y}p_{\cdot j}\\
			&=&\pause \sum_{x_i\le x}P(X=x_i)\sum_{y_j\le y}P(Y=y_j)=\pause \textcolor{red}{F_X(x)F_Y(y)}
		\end{eqnarray*}
	\end{itemize}
\end{frame}
\begin{frame}
	\frametitle{随机变量的独立性}
	\begin{itemize}[<+-|alert@+>]
		\item 二维连续型随机变量
		\begin{eqnarray*}
			F (x|y)&=&\int_{-\infty}^x\dfrac{p (u,y)}{p_Y (y)} du=P (X\le x|Y=y)\xlongequal[\mbox{不影响概率}]{\mbox{若条件事件}} P (X\le x)\\
			&=&\int_{-\infty}^xp_X(u)du\\
			&\Rightarrow& p(x,y)=p_X(x)p_Y(y)
		\end{eqnarray*}
		\item 此时，其分布函数
		\begin{eqnarray*}
			\textcolor{red}{F(x,y)}&=&P(X\le x, Y\le y)=\int_{-\infty}^x\int_{-\infty}^yp(u,v)dudv\\
			&=&\int_{-\infty}^xp_X(u)du\int_{-\infty}^yp_Y(v)dv=\textcolor{red}{F_X(x)F_Y(y)}
		\end{eqnarray*}
	\end{itemize}
\end{frame}
\begin{frame}
	\frametitle{随机变量独立性的定义}
	\begin{defi}
	 设 $n$ 维随机变量 $(X_1,X_2,\cdots,X_n)$ 的联合分布函数为 $F (x_1,
		\cdots,x_n)$, $F_i (x_i)$ 为 $X_i$ 的边缘分布函数。如果对任意的 $n$ 个实数 $x_1,x_2,\cdots, x_n$, 有
		\begin{eqnarray*}
			F(x_1,x_2,\cdots, x_n)=F_1(x_1)F_2(x_2)\cdots F_n(x_n)
		\end{eqnarray*}
		则称 $X_1,X_2,\cdots, X_n$ 相互独立。特别的，
	\end{defi}
	\pause
	\begin{enumerate}[<+-|alert@+>]
		\item 在离散随机变量场合，如果对任意 $n$ 个取值 $x_1,x_2,\cdots, x_n$, 有
		{\small     \begin{eqnarray*}
				P(X_1=x_1,X_2=x_2,\cdots,X_n=x_n)=P(X_1=x_1)P(X_2=x_2)\cdots P(X_n=x_n),
		\end{eqnarray*}}
		则称 $X_1,X_2,\cdots, X_n$ 相互独立.


		\item 在连续随机变量场合，如果对任意的 $n$ 个取值 $x_1,x_2,\cdots, x_n$, 有
		\begin{eqnarray*}
			p(x_1,x_2,\cdots, x_n)=p_1(x_1)p_2(x_2)\cdots p_n(x_n)
		\end{eqnarray*}
		则称 $X_1,X_2,\cdots, X_n$ 相互独立.
	\end{enumerate}
\end{frame}

\begin{frame}{两个随机变量独立性的另一定义}
\begin{itemize}[<+-|alert@+>]
\item 事件独立性的定义: $P(AB)=P(A)P(B)$;
\item 事件族 $\mathcal{A}$ 和 $\mathcal{B}$ 的独立性: 若对任意的 $A\in \mathcal{A}$ 和 $B\in \mathcal{B}$ 都有 $$P(AB)=P(A)P(B);$$
\item 随机变量 $X$ 和 $Y$ 的独立性:
\begin{itemize}[<+-|alert@+>]
\item 与随机变量 $X$, $Y$ 相关的事件族为:
$$\sigma(X)=X^{-1}(\mathcal{B}), \quad \sigma(Y)=Y^{-1}(\mathcal{B}), \quad \mathcal{B}:=\mathcal{B}(\mathbb{R});$$
\item 称 $X$ 和 $Y$ 独立，若它们生成的 $\sigma$代数 $\sigma(X)$ 和 $\sigma(Y)$ 独立，
\end{itemize}

\end{itemize}

\pause
\begin{defi}
称概率空间 $(\Omega, \mathcal{F}, P)$ 上的随机变量 $X$ 和 $Y$ 独立, 如果它们生成的 $\sigma$代数 $\sigma(X)=X^{-1}(\mathcal{B})$ 和 $\sigma(Y)=Y^{-1}(\mathcal{B})$ 独立，即对任意的 $A, B \in \mathcal{B}$ 都有
\[
P(X \in A, Y \in B)=P(X \in A) P(Y \in B)
\]
\end{defi}

\end{frame}


\begin{frame}{随机变量独立性的一些重要性质}
\begin{thm}
设随机变量 $X, Y$ 相互独立, 并令 $F_{X}$ 和 $F_{Y}$ 为随机向量 $(X, Y)$ 的边缘分布函数, 则对 $A, B \in \mathcal{B}$, 只要 $P(X \in A)>0, P(Y \in B)>0$ 就有
$$
F_{Y}(y | X \in A)=F_{Y}(y), \quad F_{X}(x | Y \in B)=F_{X}(x), \quad x, y\in \mathbb{R}.$$
\end{thm}%
\vspace{-0.7cm}

\pause
\begin{thm}
 设随机向量 $(X, Y)$ 有联合分布函数 $F$ 和边缘分布函数 $F_{X}, F_{Y}$ ．则随机变量 $X, Y$ 相互独立当且仅当
$$
F(x, y)=F_{X}(x) F_{Y}(y), \quad  x, y \in \mathbb{R}.$$
\end{thm}%
\
\end{frame}
\begin{frame}{随机变量独立性的一些重要性质}

\vspace{-0.7cm}
\begin{thm} 设随机变量 $X$ 和 $Y$ 独立, 则对任何博雷尔可测函数 $f$ 和 $g$, 随机变量 $f(X)$ 和 $g(Y)$ 也相互独立.
\end{thm}%

\zheng \begin{itemize}[<+-|alert@+>]
\item 因为 \( f, g \) 是博雷尔可测函数，故 \( f(X), g(Y) \) 都是随机变量.
\item 对任意的 \( A, B \in \mathcal{B} \) 有 \( f^{-1}(A), g^{-1}(B) \in \mathcal{B} \).
\item 故由 \( X \) 和 \( Y \) 的独立性知，
\[
\begin{aligned}
P\{f(X) \in A, g(Y) \in B\}  &=P\left\{X \in f^{-1}(A), Y \in g^{-1}(B)\right\} \\
& =P\left\{X \in f^{-1}(A)\right\} P\left(Y \in g^{-1}(B)\right\}  \\
&=P\{f(X) \in A\} P\{g(Y) \in B\}
\end{aligned}
\]
\item \( f(X) \) 和 \( g(Y) \) 相互独立．
\end{itemize}

\end{frame}







\begin{frame}{随机变量独立性的一些重要性质}
\begin{thm}
设 $X_{1}$ 和 $Y_{2}$ 分别是概率空间 $\left(\Omega_{1}, \mathcal{F}_{1}, P_{1}\right)$ 和 $\left(\Omega_{2}, \mathcal{F}_{2}, P_{2}\right)$ 上的随机变量. 令 $(\Omega, \mathcal{F}, P)=\left(\Omega_{1} \times \Omega_{2}, \mathcal{F}_{1} \times \mathcal{F}_{2}, P_{\mathbf{1}} \times P_{2}\right)$, 并定义
\[
X(\omega)=X_{1}\left(\omega_{1}\right), \quad Y(\omega)=Y_{2}\left(\omega_{2}\right),  \omega=\left(\omega_{1}, \omega_{2}\right) \in \Omega .
\]
则如下性质成立：
\begin{enumerate}[<+-|alert@+>]
\item $X, Y$ 是概率空间 $(\Omega, \mathcal{F}, P)$ 上的随机变量且相互独立；
\item $X$ 与 $X_{1}$ 有相同分布，而 $Y$ 与 $Y_{2}$ 有相同分布．
\end{enumerate}
\end{thm}

% 证明 首先，注意对任何 $A, B \in \mathcal{B}$ ，有
% \[
% \left\{X_{1} \in A\right\}=X_{1}^{-1}(A) \in \mathcal{F}_{1}, \left\{Y_{2} \in B\right\}=Y_{2}^{-1}(B) \in \mathcal{F}_{2} .
% \]

% 由（3．2．6）式及乘积 $\sigma$ 代数 $\mathcal{F}=\mathcal{F}_{1} \times \mathcal{F}_{2}$ 的定义有
% \[
% \{X \in A\}=\left\{X_{1} \in A\right\} \times \Omega_{2} \in \mathcal{F}, \{Y \in B\}=\Omega_{1} \times\left\{Y_{2} \in B\right\} \in \mathcal{F} .
% \]

% 根据乘积概率测度的定义得
% \[
% P(X \in A)=P_{1}\left(X_{1} \in A\right) P_{2}\left(\Omega_{2}\right)=P_{1}\left(X_{1} \in A\right)
% \]

% 和
% \[
% P(Y \in B)=P_{1}\left(\Omega_{1}\right) P_{2}\left(Y_{2} \in B\right)=P_{2}\left(Y_{2} \in B\right)
% \]

% 所以 $X$ 与 $X_{1}$ 同分布，而 $Y$ 与 $Y_{2}$ 同分布．注意
% \[
% \{X \in A, Y \in B\}=\left\{X_{1} \in A\right\} \times\left\{Y_{2} \in B\right\} \in \mathcal{F}
% \]

% 再由乘积概率测度的定义有
% \[
% P(X \in A, Y \in B)=P_{1}\left(X_{1} \in A\right) P_{2}\left(Y_{2} \in B\right)=P(X \in A) P(Y \in B)
% \]

% 这说明 $X$ 和 $Y$ 相互独立．
\end{frame}

\begin{frame}{多个随机变量独立性的定义及性质}
\begin{defi}
称概率空间 $(\Omega, \mathcal{F}, P)$ 上的 $n$ 个随机变量 $X_{1}, X_{2}, \cdots, X_{n}$ 独立，如果它们生成的 $\sigma$ 代数 $\sigma\left(X_{1}\right), \sigma\left(X_{2}\right), \cdots, \sigma\left(X_{n}\right)$ 独立，即对于任意的 $B_{1} \in \mathcal{B}, \cdots, B_{n} \in \mathcal{B}$ 都有
\[
P\left(X_{1} \in B_{1}, \cdots, X_{n} \in B_{n}\right)=P\left(X_{1} \in B_{1}\right)  \cdots P\left(X_{n} \in B_{n}\right).%
\]%
\end{defi}%
\pause
\begin{defi}
  称概率空间 $(\Omega, \mathcal{F}, P)$ 上的一族随机变量 $X_{i}, i \in I$ 独立，是指它们生成的 $\sigma$ 代数 $\sigma\left(X_{i}\right), i \in I$ 独立．
\end{defi}

\pause
\vspace{0.5cm}
\begin{thm}
随机变量 $X_{1}, X_{2}, \cdots, X_{n}$ 相互独立当且仅当随机向量 $\left(X_{1}, X_{2}, \cdots, X_{n}\right)$的联合分布函数等于其边缘分布函数的乘积，即
\[
F\left(x_{1}, x_{2}, \cdots, x_{n}\right)=F_{1}\left(x_{1}\right) F_{2}\left(x_{2}\right) \cdots F_{n}\left(x_{n}\right),  x_{1}, x_{2}, \cdots, x_{n} \in \mathbb{R} .
\]
\end{thm}

\end{frame}
\end{document}
